% [原文 p.201]
\par\medskip\noindent\textbf{命题 5.20（续）}\quad

(5)\quad 设 $U$ 为 $X$ 上使 $\Delta$ 在近旁成为 snc 的非奇异点全体。若 $g^{-1}(U)$ 非奇异且 $\Delta'|_{g^{-1}(U)}$ 为 snc，则 $(X,\Delta)$ 为 dlt 蕴含 $(X',\Delta')$ 为 dlt。

\begin{proof}
首先必须明确：有限态射下 Weil 除子的拉回要如何定义。设 $U\subset X$ 为非奇异点全体构成的集合。$X'\setminus g^{-1}(U)$ 在 $X'$ 内的余维至少为 $2$。因此，对 $X$ 上任意的 Weil 除子 $B$，只需把它限制到 $U$，拉回到 $g^{-1}(U)$，再唯一地延拓到整个 $X'$。

记 $U'=g^{-1}(U)\subset X'$。此时 $g\colon U'\to U$ 有限，且对每个 $m$ 有
\[
m(K_{U'}+\Delta'|_{U'})=g^*\bigl(m(K_U+\Delta|_U)\bigr).
\]
设 $m(K_X+\Delta)$ 为 Cartier。此时 $g^*(m(K_X+\Delta))$ 在 $X'$ 上为 Cartier，且除余维 $2$ 的子集外，它与 $m(K_{X'}+\Delta')$ 一致。于是 $g^*(m(K_X+\Delta))=m(K_{X'}+\Delta')$，故 $m(K_{X'}+\Delta')$ 为 Cartier。同样的论证表明：若 $m(K_{X'}+\Delta')$ 为 Cartier，则 $\deg g\cdot m(K_X+\Delta)$ 也为 Cartier。

为证明其余部分，考察下述纤维积与除子。
\[
\begin{tikzcd}
e'\in E' \arrow[r,hook] \arrow[d] & Y' \arrow[r,"f'"] \arrow[d,"h"] & X' \arrow[d,"g"] \\
e\in E \arrow[r,hook] & Y \arrow[r,"f"] & X
\end{tikzcd}
\]
设 $r$（$r\le\deg h=\deg g$）是 $h$ 沿 $E'$ 的分歧指数。在 $e'$ 的近旁可以如下计算。
\[
\begin{aligned}
K_{Y'} &= f'^*(K_{X'}+\Delta')+a(E',X',\Delta')E'\\
&= f'^*g^*(K_X+\Delta)+a(E',X',\Delta')E'\\
&= h^*f^*(K_X+\Delta)+a(E',X',\Delta')E',\\
K_{Y'} &= h^*K_Y+(r-1)E'\\
&= h^*f^*(K_X+\Delta)+a(E,X,\Delta)h^*E+(r-1)E'\\
&= h^*f^*(K_X+\Delta)+\bigl(ra(E,X,\Delta)+(r-1)\bigr)E'.
\end{aligned}
\]
% [原文 p.202]
于是得到 $a(E',X',\Delta')+1=r\bigl(a(E,X,\Delta)+1\bigr)$。

这里，把 $E'$ 取为 $\Delta'$ 的任意不可约分支，即可看出：$\Delta'$ 的系数至多为 $1$（或小于 $1$）当且仅当 $\Delta$ 的系数至多为 $1$（或小于 $1$）。特别地，(4) 已归约到 (2)、(3)。若 $E'$ 为例外除子，则只要两边之一不小于 $-1$，便得到 (2) 与 (3)。其余情形由推论 2.31，两边都等于 $-\infty$。(5) 通过取 $E'\subset X'\setminus g^{-1}(U)$ 得到。另外注意：(5) 的“$\Leftarrow$”并不成立。
\end{proof}

% [Tutor Note: 有限态射下拉回 Weil 除子的关键在于 $X'\setminus g^{-1}(U)$ 的余维至少为 $2$：先在非奇异部分定义拉回，再用正规概形的延拓唯一性补上余维 $2$ 以上的部分。这也是后文反复使用的手法。]

\begin{corollaryp}{5.21}
设 $x\in X$ 是正规奇点的芽。
\begin{enumerate}[label=(\arabic*),leftmargin=*]
\item $x\in X$ 为 klt，当且仅当 $x\in X$ 是指数 $1$ 的标准奇点 $0\in Y$ 在余维 $1$ 处自由作用的循环群作商。
\item 若 $x\in X$ 为终端（或标准），则它是指数 $1$ 的终端（或标准）奇点 $0\in Y$ 在余维 $2$ 处自由作用的循环群作商。
\end{enumerate}
\end{corollaryp}

\begin{proof}
设 $x\in X$ 为 klt。此时 $K_X$ 为 $\Q$-Cartier。设 $p\colon Y\to X$ 是其指数 $1$ 覆盖（定义 5.19）。由命题 5.20(3)，$\operatorname{discrep}(Y)\ge\operatorname{discrep}(X)$。由于 $K_Y$ 为 Cartier，$\operatorname{discrep}(Y)$ 是整数。由此同时得到 (2) 与 (1) 的一个方向。(1) 的逆方向由命题 5.20(4) 导出。
\end{proof}

\par\medskip\noindent\textbf{注意}\quad (2) 的逆不成立。一般地，判定终端（或标准）奇点的商何时仍为终端（或标准），并不容易。$Y$ 为三维终端的情形在定理 5.43 中处理。

下面叙述的是关于 dlt 对局部性质的最基本结果。

\begin{theoremp}[{[Elk81, Fuj85, KMM87]}]{5.22}
设 $(X,\Delta)$ 是 dlt 对，且 $\Delta$ 有效。此时 $X$ 具有有理奇点。
\end{theoremp}

\begin{proof}
在 $X$ 为射影的情形下证明。一般情形的证明要复杂得多。只用 CM 概形的对偶的证明见 [Kol97, Sec.11]。

首先由命题 2.43，可以假定 $(X,\Delta)$ 为 klt。

设 $f\colon Y\to X$ 是 $\Delta$ 的对数奇点消解。
% [原文 p.203]
表示为
\[
K_Y\equiv f^*(K_X+\Delta)-A+B,
\]
其中 $A,B$ 是无公共分支的有效 $\Q$-除子。此时 $\Supp B\subset\Exc(f)$，$\lfloor A\rfloor=0$，且 $\Supp(A+B)$ 为 snc 除子。考虑 $\Q$-除子
\[
\lfloor B\rfloor\equiv K_Y-f^*(K_X+\Delta)+A+\{-B\}.
\]
$f^*(K_X+\Delta)$ 是数值 $f$-平凡的，$A+\{-B\}$ 是满足 $\lfloor A+\{-B\}\rfloor=0$ 的 snc 除子。因此由推论 2.68，$R^jf_*\cO_Y(\lfloor B\rfloor)=0$（$j>0$）。设 $L$ 是 $X$ 上的丰富 Cartier 除子，则得到下述交换图。
\[
\begin{tikzcd}
H^i(\cO_Y(-rf^*L)) \arrow[r] & H^i(\cO_Y(\lfloor B\rfloor-rf^*L)) \\
H^i(\cO_X(-rL)) \arrow[u] \arrow[r,phantom,"="] & H^i(\cO_X(-rL)) \arrow[u,"\beta"']
\end{tikzcd}
\]
注意 Leray 谱序列
\[
H^i\bigl(X,\cO_X(-rL)\otimes R^jf_*\cO_Y(\lfloor B\rfloor)\bigr)\Rightarrow H^{i+j}\bigl(Y,\cO_Y(\lfloor B\rfloor-rf^*L)\bigr).
\]
由于 $f_*\cO_Y(\lfloor B\rfloor)=\cO_X$，且 $R^jf_*\cO_Y(\lfloor B\rfloor)=0$（$j>0$），$\beta$ 为同构。

再由定理 2.64，$H^i(\cO_Y(-rf^*L))=0$（$i<n$，$0<r$）成立。因此 $H^i(\cO_X(-rL))=0$（$i<n$）成立，由推论 5.72，$X$ 为 CM。

取 $i=n$，得到单射
\[
H^n(\cO_X(-rL))\hookrightarrow H^n(\cO_Y(-rf^*L)).
\]
因 $X$ 为 CM，用对偶（定理 5.71）得到满射
\[
H^0(Y,\omega_Y(rf^*L))=H^0(X,f_*\omega_Y\otimes\cO_X(rL))\to H^0(X,\omega_X\otimes\cO_X(rL)).
\]
于是 $f_*\omega_Y\to\omega_X$ 为满射，故 $X$ 具有有理奇点（定理 5.10）。
\end{proof}

% [Tutor Note: 分解 $K_Y\equiv f^*(K_X+\Delta)-A+B$ 是证明的技术核心：$A,B$ 分别收集系数为正、为负的部分，取下整与取负小数部分后得到 snc 除子，才能对 $\lfloor B\rfloor$ 用消没定理。]

\par\medskip\noindent\textbf{注意}\quad 定理 5.22 已发表的证明中，有两个格外简短的。一个在 [Fle81, p.36]，另一个在 [MP97, p.141]。遗憾的是，著者未能理解这些证明。

% [原文 p.204]
\begin{example}{5.23}
维数至少为 $2$ 的 Abel 簇上的锥，其顶点为 lc，但并非 CM。
\end{example}

\begin{corollaryp}{5.24}
设 $X$ 是使 $K_X$ 为 Cartier 的正规簇。此时，$X$ 具有有理奇点当且仅当 $(X,0)$ 是标准的。
\end{corollaryp}

\begin{proof}
若 $(X,0)$ 是标准的，则由定理 5.22，$X$ 具有有理奇点。反之，假定 $X$ 具有有理奇点。设 $f\colon Y\to X$ 是任意奇点消解。表示为 $K_Y=f^*K_X+E$，其中 $E$ 为 $f$-例外的。由于 $f_*\cO_Y(K_Y)=\cO_X(K_X)$，$E$ 是有效的。这意味着在 $X$ 上方的任何例外除子的差异系数都不小于 $0$。也就是说，$(X,0)$ 是标准的。
\end{proof}

在若干情形下，可以把 $\Q$-Cartier Weil 除子当作 Cartier 除子一样对待。

\begin{corollaryp}{5.25}
设 $(X,\Delta)$ 是 dlt 对，$D$ 是 $X$ 上的 $\Q$-Cartier Weil 除子。此时，$\cO_X(D)$ 是 CM 层。若 $D$ 有效，则 $\cO_D$ 为 CM。
\end{corollaryp}

\begin{proof}
由命题 2.43，不妨假定 $(X,\Delta)$ 为 klt。取 $m>0$ 使 $mD$ 为 Cartier，并取 Cartier 除子 $L$，使 $mL-mD$ 充分丰富。设 $s\in H^0(X,\cO(mL-mD))$ 为一般截面，$E$ 为其零集，$p\colon X'\to X$ 为对应的循环覆盖。

由引理 5.17(2)，$\bigl(X,\Delta+(1-\frac{1}{m})E\bigr)$ 为 klt。由命题 5.20(4)，$(X',p^*\Delta)$ 也为 klt。因此由定理 5.22，$\cO_{X'}$ 为 CM。由于 $\cO_X(D-L)$ 是 $p_*\cO_{X'}$ 的直和成分，它也为 CM（命题 5.4）。又 $\cO_X(D)$ 与 $\cO_X(D-L)$ 局部同构，故 $\cO_X(D)$ 为 CM。

设 $D$ 有效，记 $D':=p^*D$。此时 $D'$ 是 $X'$ 上的 Cartier 除子，故 $\cO_{D'}$ 为 CM。取分解 $\frac{1}{\deg p}\operatorname{Trace}_{X'/X}\colon p_*\cO_{X'}\to\cO_X$ 与 $\cO_D$ 的张量积，便得到 $\cO_D\to p_*\cO_{D'}$ 的分裂。于是 $\cO_D$ 为 CM。
\end{proof}

% [Tutor Note: 后一段原文把 $D'$ 所在的簇写作 $X$，按上下文应为 $X'$。$D'=p^*D$ 是 $X'$ 上的 Cartier 除子，$\cO_{D'}$ 也在 $X'$ 上。底本如此，照文意译出。]

\begin{propositionp}{5.26}
设 $(X,\Delta)$ 是射影的 dlt 对，$S,D$ 是 $\Q$-Cartier Weil 除子。设 $S$ 有效，并进一步假定 $D$ 在 $X$ 上余维 $2$ 处为 Cartier。此时，由限制得到的列
% [原文 p.205]
\[
0\to\cO_X(D-S)\to\cO_X(D)\to\cO_S(D|_S)\to0
\]
是正合的。其中 $\cO_S(D|_S)$ 是 $S$ 上满足下述条件的唯一 $S_2$ 凝聚层。设 $U\subset X$ 是使 $D$ 成为 Cartier 的开集之中的极大者，则 $\cO_S(D|_S)|_{U\cap S}=\cO_{U\cap S}(D|_{U\cap S})$。（当然，若进一步 $S$ 正规，则这与已有的定义一致。）
\end{propositionp}

\begin{proof}
这样的 $\cO_S(D|_S)$ 唯一，由推论 5.25 知 $\cO_S$ 为 CM，故这一点是显然的。设 $p\colon X'\to X$ 如前证明，记 $S':=p^*S$。$D':=p^*D$ 为 Cartier，故下述自然列
\[
0\to\cO_{X'}(D'-S')\to\cO_{X'}(D')\to\cO_{S'}(D'|_{S'})\to0
\]
是正合的。取此正合列在 $X$ 上的像，并取出适当的直和成分，得到
\[
0\to\cO_X(D-S-L)\to\cO_X(D-L)\to F\to0,
\]
其中 $F$ 是 $p_*\cO_{S'}(D'|_{S'})$ 的直和成分。只需证明 $F=\cO_S(D-L|_S)$。

问题在于，从所用 $X'\to X$ 的构造并不能直接看出 $p_*\cO_{S'}=p_*\cO_{X'}\otimes\cO_S$。然而，在使 $p$ 平展的开集（$X$ 的）上等式成立。因此 $F|_{U\setminus E}\cong\cO_S(D-L|_S)|_{U\setminus E}$。由于 $D'$ 为 Cartier，由推论 5.25，$p_*\cO_{S'}(D'|_{S'})$ 为 CM。又由 $\cO_S(D)$ 的定义，$\cO_S(D-L|_S)$ 为 $S_2$，且 $S\cap(X\setminus U)$ 在 $S$ 内的余维至少为 $2$。因此 $F|_{X\setminus E}\cong\cO_S(D-L|_S)|_{X\setminus E}$。

由于可以让 $E$ 变动，即得 $F\cong\cO_S(D-L|_S)$。
\end{proof}

\begin{corollaryp}{5.27}
（dlt 对的 Serre 对偶）设 $(X,\Delta)$ 是特征 $0$ 的域上的纯 $n$ 维射影 dlt 对，$D$ 是 $X$ 上的 $\Q$-Cartier Weil 除子。此时，$H^i(X,\cO_X(D))$ 与 $H^{n-i}(X,\omega_X(-D))$ 互为对偶。
\end{corollaryp}

\begin{proof}
这由推论 5.25 与定理 5.71 得到。
\end{proof}

最后，不加证明地叙述标准奇点的一个有用性质。

\begin{theorem*}[{[Kol83]}]{5.28}
设 $f\colon X\to S$ 是平坦态射，纤维只有标准奇点。此时下列成立。
\begin{enumerate}[label=(\arabic*),leftmargin=*]
\item $\omega^{[q]}_{X/S}$（即 $\omega^{\otimes q}_{X/S}$ 的双对偶）的构造，是与基变换可交换的操作。
\item 若 $f$ 固有，则 $s\mapsto\chi(X_s,\omega^{[q]}_{X_s})$ 是 $S$ 上的局部常值函数。
\end{enumerate}
\end{theorem*}

% [原文 p.206]
\section{三维终端奇点与标准奇点}
\takeaway{本节研究三维终端奇点与标准奇点。核心技术是结构定理 5.34、5.35：指数 $1$ 的标准奇点，其一般超曲面截面为 Du Val 奇点，且存在差异系数为 $1$ 的例外除子。完整分类（定理 5.43）不涉及。}

本节的目标是详细研究三维终端奇点与标准奇点。指数 $1$ 的三维标准奇点与终端奇点的结构定理，将在定理 5.34 与定理 5.35 中证明。任意三维终端奇点的分类定理 5.43 以此结果为依据。这一详细分类有时以“终端引理”之名被引用，但本书不需要它。

\begin{notation}{5.29}
本节的讨论在代数、解析两种设定下都有效。因此，设 $(0\in X)$ 是给定了闭点 $0$ 的 $\C$ 上的正规仿射簇，或正规 Stein 簇。注意：$0$ 不一定是 $X$ 的孤立奇点。

这里把 $X$ 看作 $0$ 的近旁。因此，证明过程中如有必要，可以随时把 $X$ 换为更小的近旁，这一点不再声明。\hfill$\square$
\end{notation}

% [Tutor Note: 把 $X$ 换为任意小近旁的做法，在解析设定下对应取足够小的 Stein 邻域。所有陈述都是芽层面的，故这样做不影响结论，但读者须留意后文出现的 $X$ 都指某个够小的近旁。]

下一个结果属于 [Rei80]，它是通向分类的第一步。

\begin{lemmap}{5.30}
设 $(0\in X)$ 是指数 $1$ 的三维标准奇点，$0\in H\subset X$ 是过 $0$ 的一般超曲面截面。此时，$(0\in H)$ 是 Du Val 奇点，或椭圆奇点。
\end{lemmap}

\begin{proof}
由定理 5.22，$X$ 为 CM，故 $H$ 同样为 CM。由于 $H$ 是一般超曲面截面，$H$ 在余维 $1$ 处非奇异。于是 $H$ 是正规曲面。设 $g\colon H'\to H$ 是任意奇点消解。此时须确认：或者 $g_*\omega_{H'}=\omega_H$（即 $H$ 为 Du Val），或者 $g_*\omega_{H'}=\mathfrak{m}_{0,H}\omega_H$（即 $H$ 为椭圆的）。

设 $B_0X\to X$ 是极大理想 $\mathfrak{m}_{0,X}\subset\cO_X$ 的爆破，$f\colon Y\to X$ 是支配 $B_0X$ 的奇点消解。
% [原文 p.207]
$f^*\mathfrak{m}_{0,X}\subset\cO_Y$ 是定义 Cartier 除子 $E\subset Y$ 的理想层。设 $H'=f^{-1}H\subset Y$ 是 $H$ 在 $f$ 下的双有理变换，则 $H'+E=f^*H$。由于线性系 $|H'|$ 无基点，$H'$ 非奇异。由于 $(0\in X)$ 为标准的，$\omega_Y=f^*\omega_X(F)$ 成立，其中 $F$ 为有效除子。于是由伴随，
\[
\begin{aligned}
\omega_{H'} &= \omega_Y(H')|_{H'}\\
&= f^*\bigl(\omega_X(H)\bigr)(F-E)|_{H'}=f^*\omega_H\bigl((F-E)|_{H'}\bigr)
\end{aligned}
\]
成立。因此
\[
\begin{aligned}
f_*\omega_{H'} &\supset \omega_H\otimes f_*\cO_{H'}(-E|_{H'})\\
&= \omega_H\otimes\bigl(f_*\cO_Y(-E)|_H\bigr)=\mathfrak{m}_{0,H}\cdot\omega_H.
\end{aligned}
\]
由此，$(0\in H)$ 为 Du Val 或椭圆的。
\end{proof}

\begin{remark}{5.31}
上述证明在所有维数下都有效。由此可以证明：指数 $1$ 的标准奇点 $0\in X$ 的过 $0$ 的一般超曲面截面 $0\in H\subset X$，具有有理奇点或椭圆奇点。
\end{remark}
% [Tutor Note: 原书此处印作“注意 5.31”，与全书 remark 的标签“注”不同。为与其余各章统一，本书一律排作“注”，所指相同。]

\begin{definition}{5.32}
设 $(0\in X)$ 是三维奇点。若它的一般超曲面截面 $0\in H\subset X$ 为 Du Val 奇点，则称 $(0\in X)$ 为复合 Du Val（compound Du Val）奇点或 cDV 奇点。
\end{definition}

\begin{definition}{5.33}
设 $(0\in X)$ 是复解析奇点。记 $\operatorname{Aut}(0\in X)$ 为所有自同构 $\varphi\colon(0\in U_1)\cong(0\in U_2)$ 的集合，其中 $0\in U_i\subset X$ 是适当的欧氏开邻域。若这样的两个同构 $\varphi_1,\varphi_2$ 在 $0$ 的某个邻域上一致，则视二者为同一。在 $\operatorname{Aut}(0\in X)$ 上引入拓扑是容易的，但此处只把 $\operatorname{Aut}(0\in X)$ 当作离散群来处理。
\end{definition}

引理 5.30 的两种情形表现出完全不同的行为。两者都有若干已知的刻画。

\begin{theorem*}{5.34}
设 $(0\in X)$ 是三维奇点，$0\in H\subset X$ 是过 $0$ 的一般超曲面截面。此时，下列条件等价。
\begin{enumerate}[label=(\arabic*),leftmargin=*]
\item $0\in H$ 为 Du Val 奇点。
\item $(X,H)$ 是标准的对。
% [原文 p.208]
\item $X$ 是指数 $1$ 的标准奇点，且对任意奇点消解 $f\colon Y\to X$ 与任意例外除子 $E\subset f^{-1}(0)$，有 $a(E,X)\ge1$。
\end{enumerate}
\end{theorem*}

\begin{theorem*}{5.35}
设 $0\in X$ 是指数 $1$ 的三维标准奇点。此时，下列条件等价。
\begin{enumerate}[label=(\arabic*),leftmargin=*]
\item 一般超曲面截面 $0\in H\subset X$ 具有椭圆奇点。
\item 在射影双有理态射 $f\colon Y\to X$ 之中，存在这样的：$f^{-1}(0)$ 非空，为纯余维 $1$ 的集合；$f\colon Y\setminus f^{-1}(0)\to X\setminus\{0\}$ 为同构；且 $K_Y=f^*K_X$。并且可以选取 $Y$，使 $X$ 上的 $\operatorname{Aut}(0\in X)$-作用全部延拓为 $Y$ 上的作用。
\item 对任意奇点消解 $f\colon Y\to X$，存在满足 $a(E,X)=0$ 的除子 $E\subset f^{-1}(0)$。
\end{enumerate}
\end{theorem*}

% [Tutor Note: 定理 5.35(2) 要求存在一个“极小”的双有理态射：例外集为纯余维 $1$ 的椭圆构型，且 $K_Y$ 为 $f$-平凡。这样的 $Y\to X$ 在三维标准奇点的收缩理论中扮演中心角色。]

\par\medskip\noindent\textbf{5.36}（定理 5.34 与定理 5.35 的证明方针）\quad 如下同时证明这两个定理。

首先证明定理 5.34(1) $\Rightarrow$ 定理 5.34(2) $\Rightarrow$ 定理 5.34(3)。容易看出定理 5.35(2) $\Rightarrow$ 定理 5.35(3)。然后证明定理 5.35(1) $\Rightarrow$ 定理 5.35(2)。

由引理 5.30，定理 5.34(1) $\Rightarrow$ 定理 5.34(3) 与定理 5.35(3) $\Rightarrow$ 定理 5.35(1) 等价；定理 5.34(3) $\Rightarrow$ 定理 5.34(1) 与定理 5.35(1) $\Rightarrow$ 定理 5.35(3) 等价，这一点也显然。于是两个定理的证明同时完成。

下面开始证明。

从定理 5.34(1) $\Rightarrow$ 定理 5.34(2) 开始。

由于 $H$ 的嵌入维数至多为 $3$，可知 $X$ 的嵌入维数至多为 $4$。于是 $0\in X$ 是超曲面奇点，特别为 CM，且 $\omega_X$ 局部自由。因此由命题 5.73，得到正合列
\[
0\to\omega_X\to\omega_X(H)\to\omega_H\to0.
\]
设 $f\colon Y\to X$ 是 $(X,H)$ 的对数奇点消解，$H'=f^{-1}_*H$。在 $Y$ 上，正合列
% [原文 p.209]
\[
0\to\omega_Y\to\omega_Y(H')\to\omega_{H'}\to0
\]
存在。对它作用 $f_*$。由于 $R^1f_*\omega_Y=0$（参见推论 2.68），得到水平方向正合的交换图
\begin{equation}
\begin{tikzcd}
0 \arrow[r] & f_*\omega_Y \arrow[r] \arrow[d,"\delta"'] & f_*\omega_Y(H') \arrow[r] \arrow[d,"\beta"] & f_*\omega_{H'} \arrow[r] \arrow[d,"\gamma"] & 0\\
0 \arrow[r] & \omega_X \arrow[r] & \omega_X(H) \arrow[r,"\alpha"'] & \omega_H \arrow[r] & 0
\end{tikzcd}
\tag{5.3}
\end{equation}
由于 $H$ 是标准的，$\gamma$ 为同构。由中山引理，
\[
f_*\omega_Y(H')=\omega_X(H)
\]
成立。与记号 2.26 同样地表示为
\[
\omega_Y(H')=f^*\omega_X(H)\otimes\cO_Y\bigl(\textstyle\sum a_iE_i\bigr),
\]
则 $\sum a_iE_i$ 有效当且仅当 $(X,H)$ 是标准的。对上式作用 $f_*$，得
\[
f_*\omega_Y(H')=\omega_X(H)\otimes f_*\cO_Y\bigl(\textstyle\sum a_iE_i\bigr).
\]
把它与上式比较，即可断定 $\sum a_iE_i$ 有效，且 $(X,H)$ 是标准的。

其次，假定定理 5.34(2) 成立。沿用上述记号，记 $f^*H=H'+\sum b_iE_i$。由假定，$K_X+H$ 为 $\Q$-Cartier，故 $K_X$ 同样如此。于是可以表示为
\[
K_Y\equiv f^*K_X+\textstyle\sum (a_i+b_i)E_i.
\]
这表明 $(X,\varnothing)$ 是标准的。并且，若 $f(E_i)\subset H$，则 $a(E_i,X)=a_i+b_i\ge b_i$ 为正。为得到定理 5.34(3)，必须证明 $K_X$ 为 Cartier。

把 $K_Y+H'\equiv f^*(K_X+H)+\sum a_iE_i$ 限制到 $H'$，得到 $K_{H'}\equiv f^*K_H+\sum a_i(E_i|_{H'})$。由定理 4.20，这表明 $H$ 为 Du Val。正如证明开头所见，$X$ 是超曲面奇点，特别地 $K_X$ 为 Cartier。

假定定理 5.35(2) 成立。此时，存在位于 $X$ 上方的除子 $E$，满足 $a(E,X)=0$ 与 $\operatorname{center}_X(E)=\{0\}$。由推论 2.32(1)，这样的除子出现在所有奇点消解中。于是得到定理 5.35(3)。

% [Tutor Note: 5.34(1)$\Rightarrow$(2) 的论证只用了 $H$ 标准与 $\omega_X(H)$ 的生成性，最后由中山引理把 $f_*\omega_Y(H')=\omega_X(H)$ 换成 $\sum a_iE_i$ 的有效性。]

% [原文 p.210]
其次，证明定理 5.35(1) $\Rightarrow$ 定理 5.35(2)。

设 $0\in H\subset X$ 为一般超曲面截面。由假定，$(0\in H)$ 是具有不变量 $k$（$k\ge1$）的椭圆奇点（定理 4.57）。在定理 4.57 中，对每个椭圆奇点构造了具体的双有理态射 $B^w_0H\to H$。这里的目标是：构造与各一般超曲面截面 $H\subset X$ 所定的态射 $B^w_0H\to H$ 相容的态射 $f\colon Y\to X$。如定理 4.57 所述，必须分别处理 $3$ 种情形。眼下先假定事情顺利进行，已到达下述状态。

\textbf{假定}：存在具有下列性质的双有理态射 $f\colon Y\to X$。
\begin{enumerate}[label=(\arabic*),leftmargin=*]
\item $f^{-1}$ 在 $\{0\}$ 之外为同构。
\item 设 $0\in H\subset X$ 是过 $0$ 的一般超曲面截面，表示为 $f^*H=H'+E$，其中 $E$ 为 $f$-例外的。此时，$f\colon H'\to H$ 与定理 4.57 所描述的加权爆破同构，且 $\cO_Y(H')$ 为 $f$-丰富。
\item $X$ 上的 $\operatorname{Aut}(0\in X)$-作用延拓为 $Y$ 上的作用。
\end{enumerate}

假定 $f\colon Y\to X$ 存在，说明如何完成证明。

由定理 4.57，$H'$ 只有 Du Val 奇点，且 $K_{H'}\equiv-E|_{H'}$ 成立。于是 $Y$ 在 $H'$ 的近旁正规。由于 $H'$ 为 $f$-丰富，$H'$ 的任意近旁都包含形如 $E\setminus\{\text{有限点集}\}$ 的部分集合。把 $Y$ 换为其正规化。注意这不会改变 $H'$ 的近旁。（更仔细地讨论可知 $Y$ 自身正规，但这里不需要。）

存在适当的有效例外除子 $F$，使 $K_Y\equiv f^*K_X+F$。取伴随，得到下式。
\[
\begin{aligned}
-E|_{H'}\equiv K_{H'} &= (K_Y+H')|_{H'}=\bigl(f^*(K_X+H)+F-E\bigr)|_{H'}\\
&\equiv (F-E)|_{H'}.
\end{aligned}
\]
特别地，$F|_{H'}\equiv0$。$H'$ 为 $f$-丰富，且与所有例外除子相交。因此若 $F|_{H'}\equiv0$，则 $F$ 有效且 $F\subset f^{-1}(0)$，由此得到 $F=0$（引理 3.39）。这意味着定理 5.35(2)。

还须证明满足上述条件的态射 $f\colon Y\to X$ 存在。设 $0\in H\subset X$ 是具有不变量 $k=k(H)$ 的一般超曲面截面。
